\section{Zusammenfassung}
\label{sec:zusammenfassung}

Die Arbeit untersucht die Zahnfußtragfähigkeit der Prüfverzahnung kst-E in der Paarung Stahl gegen PA46-GF30 mit Berechnung und Laufversuch. Die Geometrie der Verzahnung wurde nach DIN ISO 21771 in einer eigenen Kette erzeugt und gegen STplus abgeglichen; aus ihr entstanden das Sektormodell mit fünf Zähnen, die Fußrundungsvarianten und die Eingaben der analytischen Verfahren. Das FE-Modell übernimmt die Faserorientierung der Spritzgießsimulation über die CONVERSE-Zuordnung und rechnet vier Werkstoffstufen vom isotropen linearen bis zum anisotropen elastisch-plastischen Modell an 99 Stellungen über dem Wälzweg bei $8$, $12$ und $16$\,Nm. Die Voruntersuchungen legten Netz, Elementtyp, Kopfkante und Körper fest und vermaßen die Prüfräder vor dem Versuch.

Die Spritzgießsimulation trifft die Struktur des CT-Befunds aus FVA 856 I im Fußgrund, im Kern und im örtlichen Orientierungsgrad der Fußrundungen, nicht aber die Faserrichtung in oberer Flanke und Kopf. Über den Umfang des Rades streut der Anisotropie-Index in Wellen von $12$ bis $15$ Zähnen. Die Anisotropie erhöht die größte Zugspannung der Fußrundung gegenüber der isotropen linearen Stufe um $15$ bis $20$\,\%, die Plastizität senkt sie um $32$ bis $41$\,\%; der Orientierungseffekt zwischen Zuordnung und CT-Zonen liegt bei $-8$ bis $+5$\,\%. Der Nachweis nach Stommel et al. besteht mit der plastischen isotropen Stufe den statischen Nachweis und sagt Lastwechselzahlen von $3{,}5 \cdot 10^{5}$ bis $1{,}7 \cdot 10^{3}$ bis zum ersten Anriss vorher. Die eigene Erweiterung mit der Hill-Vergleichs\-spannung verschiebt den kritischen Knoten von der Stirnseite in die Zahnmitte; ihre Höhe hängt von der Klassenwahl ab. S-Life Plastics bestätigt Spannungen, statischen Nachweis und Amplitude; der Unterschied in der zyklischen Auslastung ist die Kappe der zulässigen Amplitude. Die Patentkonturen der Fußrundung senken die Zugspannung an C um $25$ bis $28$\,\%, unter der plastischen Karte um $16$ bis $21$\,\%.

Neun Laufversuche auf drei Stufen ergaben eine Bruchlinie mit $\lg N_{\mathrm{L},50\%} = 8{,}721 - 0{,}1617 \, T_{2}$ und einer Standardabweichung von $0{,}13$; die Werte für $10$\,\% nach DIN 50100 und VDI 2736 Blatt 4 liegen beim $0{,}68$- bis $0{,}74$-Fachen. Die ersten Anrisse wurden bei $8$\,Nm zwischen $9$ und $31$\,\% der Laufzeit sichtbar. Gegenüber POM auf demselben Prüfstand trägt PA46-GF30 bei $16$\,Nm das $1{,}6$- bis $2{,}1$-Fache an Lastwechseln, gegenüber PEEK weniger. Die Vorhersage des Nachweises liegt um die Faktoren $75$ bis $820$ unter der Bruchlinie und um mindestens den Faktor $9$ unter dem ersten Anriss; sie ist auf der sicheren Seite, und der Abstand wächst mit dem Moment. Mit den Spannungen der plastischen Stufen folgt aus den Versuchen eine Steigung von $k = 10$ bis $12$ nahe der synthetischen Linie, so dass eine am Versuch kalibrierte Linie vor allem die Lage anheben würde. \Note{PH: Einordnung von VDI 2736, RIKOR und FVA-Workbench, sobald die Werte vorliegen.}

\section{Ausblick}
\label{sec:ausblick}

Für die weiteren Laufversuche an dieser Verzahnung folgt aus den Anrissfenstern eine Prüfanweisung. In der ersten Million Lastwechsel ist etwa stündlich, also alle $200\,000$ Lastwechsel, anzuhalten und die Fußrundung der markierten Zähne zu fotografieren; zeigt die erste Kontrolle schon Anrisse, ist beim nächsten Rad nach $100\,000$ oder $50\,000$ Lastwechseln zu prüfen. Jeder Halt wird mit Lastwechselzahl, Moment vor und nach dem Nachspannen, Öltemperaturen und Befund je Zahn dokumentiert, und der erste sichtbare Anriss wird ausdrücklich festgehalten. Das Startmoment wird um einen festen Anteil über dem Sollwert eingestellt, damit der Sollwert in der Mitte des Nachspannintervalls liegt. Ein Datalog des Analogsignals der Messwelle würde den Momentverlust kontinuierlich erfassen; die Schwankung der Anzeige bei $2250$\,1/min lässt sich mit zwei Drehzahlen in ihrer Nähe einer Erfassungsursache oder einer Resonanz zuordnen. Mit dokumentierten Anrissen wird die zweite Stufe des Vergleichs auf allen Horizonten möglich, und die lokale Wöhlerlinie am kritischen Knoten kann als Anrisslinie kalibriert werden.

Auf der Modellseite steht das Abwälzen in einer Simulation mit Vorbelastung und Reibung aus, das die Lastgeschichte eines Knotens als Vorgang statt als Folge von Stellungen liefert. In dieser Arbeit wurde nach dem Buch jede Stellung einzeln gerechnet, weil die Werkstoffstufen und die Kantenfrage Vorrang hatten. Ein Schliffbild eines kst-E-Rades in der Stirnebene entscheidet, ob die Fasern in oberer Flanke und Kopf wie in der Simulation entlang der Zahnbreite oder wie im CT in der Ebene liegen; davon hängt ab, welche Orientierung die Stufen W3 und W4 tragen sollen. Die Wahl des Zahns für die Zuordnung ist als Eingabe zu führen, oder die Rechnung läuft über mehrere Zähne des Vollrades. Die Klasse der Abminderungsfaktoren unter der Hill-Vergleichs\-spannung ist mit Versuchen an gekerbten Proben oder am Zahn festzulegen. Die Fußrundungsvarianten sind mit den Nachtläufen der anisotropen Stufen auf die Kombination aus Geometrie und Fertigung zu prüfen und als gefertigte Räder im Laufversuch zu bestätigen; die Varianten V1 bis V4 der Messspur und ein Pulsatorversuch an denselben Rädern ergänzen das Bild. Eine zweite Drehzahl trennt den Einfluss der Dehnrate von dem der Temperatur.

Die Werkzeugkette aus \secref{sec:werkzeugkette} beschreibt mit ihren Datenkontrakten jede Größe mit Normsymbol, Einheit und Quelle, von der Geometrie nach Norm über die Eingabedateien des FE-Modells bis zur Auswertung. Dieselben Kontrakte eignen sich für die Messprotokolle und Rohdaten des Labors, von der Arbeitsmappe des Prüfstands bis zu den Konturscans. Eine Masterarbeit kann damit auf einer durchgängigen Kette vom Versuch bis zur Bemessung aufsetzen.

% Herkunft:
% 8.1: Kapitel 3 bis 7 dieser Arbeit (Zahlen aus kap5_1, kap5_2, kap5_3, kap5_4, kap6_3, kap6_4, kap7_1).
% 8.2 Absatz 1 (Prüfanweisung): Nutzer 2026-10-18 (stündlich/200 000 in der ersten Million, 100 000 oder 50 000
%   beim nächsten Rad, erster sichtbarer Anriss dokumentieren; Startmoment über dem Sollwert; Datalog aus Spalte J
%   der Mappe; Drehzahlvariation gegen die Schwankung bei 2250 1/min aus kap6_1); ADR-118 Amendment 2026-10-18
%   Abend Punkt 3 und 5; E8 (Roadmap).
% 8.2 Absatz 2: Plan 2026-10-14 Kapitel 8 (Abwälzen S9, M7-Varianten V1 bis V4, Schliffbild, Pulsator, zweite
%   Drehzahl); ADR-119 (Zahnwahl als Eingabe, Vollrad); ADR-118 E7 (Klasse); kap5_4 (Nachtläufe der Varianten).
% 8.2 Absatz 3: Plan 2026-10-14 Schreibregel 2 (Programmierunterstützung nur einmal im Methodenteil 3.5;
%   pydantic-Kontrakte als Brücke zur Masterarbeit; hier ohne Nennung des Werkzeugs).
% Bildwunsch: keiner. Bibwunsch: keiner. Nomenklatur: keine neuen Zeichen.
% Offen: VDI/RIKOR/Workbench in der Zusammenfassung.
